% ============================================================
% ABSTRACT
% ============================================================

\chapter*{ABSTRACT}
\addcontentsline{toc}{chapter}{Abstract}

As the number of mobile applications continues to grow and network usage demands diversify, controlling Quality of Service (QoS) on mobile devices has become a critical technical challenge. Existing solutions on Android either require Root privileges --- impractical for most users --- or only provide basic firewall functionality without priority-based bandwidth scheduling capabilities.

This thesis presents the design and implementation of \textbf{QoS Scheduler} --- an Android application capable of intervening at the network layer (Layer 3) to perform per-application bandwidth scheduling \textbf{without requiring Root access}. The system utilizes the TUN virtual network interface through the Android VpnService API to intercept, analyze, and schedule all IP traffic on the device.

The main contributions of this thesis include:
\begin{itemize}
    \item A Data Plane module capable of bit-level IPv4/IPv6 and TCP/UDP header parsing, including IPv6 Extension Header Chaining support.
    \item Implementation of the \textbf{Token Bucket} algorithm for rate limiting and \textbf{Weighted Fair Queuing (WFQ)} for priority-weighted bandwidth allocation (HIGH/MEDIUM/LOW).
    \item A UID Resolution mechanism based on 5-tuple information for per-application traffic classification.
    \item A stable TCP/UDP Relay system compatible with multiple Android versions, including custom ROMs such as HyperOS/MIUI.
\end{itemize}

Experimental results demonstrate the system's ability to accurately limit bandwidth, allocate network resources by weight with less than 5\% deviation, and mitigate Bufferbloat for high-priority applications.

\vspace{1cm}
\noindent\textbf{Keywords:} Quality of Service, QoS, Android, VPN, Token Bucket, Weighted Fair Queuing, Packet Analysis, Bandwidth Scheduling.
